\begin{papercard}{magoulas2026}
\cf{Stages} \stagetag{D} \stagetag{R} (also \stagetag{A}). A large-scale descriptive analysis of Google Play \emph{Data Safety} disclosures and declared permissions for 49{,}578 popular apps, collected from a European terminal in May 2023. Exemplary for treating structured store disclosures as \emph{presentations}, not verified practice.

\cf{Aim and framing} Examine how developers disclose collection, sharing, security and deletion through the Data Safety section and permissions five years after the GDPR. It does \emph{not} assess GDPR compliance or audit behaviour \pc{magoulas2026}{2}{Rather than assessing GDPR compliance per se}. These disclosures are self-reported and not systematically verified \pc{magoulas2026}{2}{They are reported by the developers and they are not}.

\cf{Sampling and capture} The Top Free, Top Paid and Top Grossing lists across all categories gave 65{,}747 entries, 49{,}578 after de-duplication \pc{magoulas2026}{4}{The primary criterion for selecting these}. Collected via a European terminal \pc{magoulas2026}{5}{The download of this data was done by a}. Recorded fields: name, rating, numbers of ratings and reviews, downloads, version, in-app purchases, last update, minimum Android version, price, permissions, and all Data Safety fields \pc{magoulas2026}{4}{The data collected from Google Play included app name, rating}.

\cf{Data Safety structure} Three parts: \emph{Data Shared}, \emph{Data Collected}, \emph{Security Practices}. 14 data categories with sub-types: location, personal info (incl.\ race/ethnicity, beliefs, sexual orientation), financial, \emph{health and fitness}, messages, photos/videos, audio, files, calendar, contacts, app activity, web browsing, app info/performance, device IDs (Table 2). Security items: Families Policy, encrypted in transit / not encrypted, independent review, data can/can't be deleted, deletion on request \pc{magoulas2026}{11}{The Data Safety section is organized into three parts}.

\cf{Absence as a variable} 2\,\% had no policy link. Links were not checked for function or content \pc{magoulas2026}{6}{did not provide a link to a privacy policy}. 25.44\,\% had completed \emph{no} Data Safety part, and analyses of content used only the 36{,}964 with at least one part, an explicit denominator rule \pc{magoulas2026}{12}{had not completed any of the}. A logistic regression on non-completion found it concentrated among older and less-installed apps, with paid apps more likely to be incomplete \pc{magoulas2026}{16}{The results showed that non-completion was clearly associated with}. Non-completion is therefore \emph{not random}, a missing-data insight.

\cf{Analyses} Correlations; OLS of rating score on permission breadth with category, price and age controls, VIF checks, Breusch--Pagan heteroskedasticity test with HC3 robust SEs, and nested models to test interpretations \pc{magoulas2026}{10}{The model was estimated using ordinary least squares (OLS)}. Descriptive yearly trends. $\chi^2$ tests of category vs.\ data type, with Cramér's $V$ effect sizes and standardised residuals, stressing effect size over significance in large samples \pc{magoulas2026}{17}{Given the large sample size, statistical signifi}.

\cf{Key findings} Free apps request more permission categories than paid apps. Breadth rose about 0.24 categories per year. 11\,\% explicitly declared data \emph{not} encrypted in transit, and 34\,\% declared \emph{no} user-initiated deletion. These are explicit negative declarations, not verified properties \pc{magoulas2026}{15}{explicitly declared that data are not encrypted in transit and 34}. Permissions and Data Safety did not correlate, which reflects their different purposes \pc{magoulas2026}{16}{A correlation analysis between permission categories and Data Safety section disclo}.

\cf{Limitations stated} Metadata and self-reports only; permissions are not runtime behaviour; a European IP does not imply GDPR applicability; top lists are not the population \pc{magoulas2026}{19}{This study has several limitations that should be acknowledged}.

\cf{Use for this thesis} The best template for the \emph{store-disclosure surface} (RQ4): record the full Data Safety schema per app, distinguish \emph{not completed} from \emph{declared none} and \emph{declared negative} (e.g.\ ``data can't be deleted''), keep separate denominators, and phrase all findings as disclosed presentation. Its careful limitation wording is a model for the thesis' claims section. The health-and-fitness Data Safety category is a direct field for comparing with policies and the UI (as Malki does).
\end{papercard}
